\chapter{The Grammar of Experiments}
\label{chap:the-grammar-of-experiments}

% Closed question: How do multiple records combine without collapsing into one
% undifferentiated record?

\begin{chapteressay}

A weather station stands in a field. It has a thermometer, barometer,
anemometer, rain gauge, hygrometer, and clock. At noon, the station reports:
temperature 22 C, pressure 101.2 kPa, wind from the west at 12 knots, no rain,
humidity 48 percent.

No single instrument produced that state. The state is a composition of records.

The thermometer did not measure wind. The anemometer did not measure pressure.
The rain gauge did not measure humidity. Each instrument carried its own
admissible marks under its own calibration. Yet the station report is one
physical claim: the atmosphere here, at this time, has this measured state.

That claim only works because the records combine without dissolving into one
another. The temperature remains temperature. The pressure remains pressure.
The wind direction remains wind direction. The clock binds them to the same
time. If one instrument is miscalibrated, the state may be inconsistent. If the
clock is wrong, readings taken at different moments may be falsely combined. If
the rain gauge is clogged, the station may report a physically impossible
combination: wet ground, falling pressure, dark clouds, and no rainfall.

The same structure appears in the speedometer thread. A car's speedometer,
odometer, and clock form an experimental state. If the speedometer reads 60 mph
for one hour, the odometer should advance by about 60 miles. If the odometer
advances by 40 miles, the disagreement is not a new kind of motion. It is a
diagnosis: one record is wrong, delayed, miscalibrated, or being interpreted
outside its conditions.

This chapter names the grammar that lets multiple records become one
experimental state. The grammar must be strict enough to detect disagreement
and flexible enough to keep distinct records distinct. It must not collapse all
records into a single undifferentiated number. It must not treat incompatible
records as if they simply averaged into truth. It must not force unresolved
states to close before the instrument has earned closure.

The closed question is:

\begin{center}
\emph{How do multiple records combine without collapsing into one
undifferentiated record?}
\end{center}

The answer is compatibility. Records combine through a protocol that names
their source, calibration, time, role, and admissible relation to the other
records. A combined experiment is not a heap of readings. It is a structured
state.

This structure will matter everywhere later. A classical mechanical state
combines position, momentum, time, and calibration. A quantum experiment may
hold incompatible alternatives open until detection. A relativistic measurement
combines local clocks through a reference geometry. A thermodynamic record
combines energy, entropy, memory, and closure. Each later chapter depends on
the grammar introduced here: records can compose only when their source
identities are preserved.

\end{chapteressay}

% Phenomenon arcs use: 1/3 setup -> phenomenon box -> 2/3 structural interpretation.

\section{Experimental State}
\label{se:experimental-state}

An experimental state is a compatible bundle of records. It is not a single
reading, and it is not a pile of unrelated readings.

The weather station gives the clean example. Temperature, pressure, wind,
humidity, rainfall, and time are distinct records. Each has its own instrument
and calibration. The station state exists only when those records are tied to
the same site and the same reporting interval. A temperature reading from noon
and a wind reading from midnight do not form a noon weather state, even if both
readings are individually accurate.

The speedometer version is the tuple of speedometer, odometer, and clock. Speed
is a ratio of distance to time. The odometer carries accumulated distance. The
clock carries elapsed time. The three records constrain one another. A stable
60 mph speed reading for one hour should agree with 60 miles of odometer
advance. If it does not, the combined state fails compatibility.

Compatibility is stronger than coexistence. Two readings can coexist on a page
without forming a state. To form a state, the records must have a named relation
that the experiment is prepared to check. The weather station says these
readings describe the same atmosphere. The car dashboard says these readings
describe the same trip. A laboratory report says these measurements describe
the same sample under the same protocol.

This is why experimental metadata is not clerical fluff. Time, location,
instrument identity, calibration date, sample ID, and procedure are part of the
state. Without them, records cannot be combined honestly. A pressure reading
without altitude, a temperature reading without sensor placement, or a
speedometer reading without tire calibration may still be a mark, but it is not
yet a reliable component of an experimental state.

The grammar of experiments begins with the insistence that composition preserve
source. Temperature does not become pressure because both appear in the same
report. Speed does not become distance because speed and odometer readings
constrain each other. A record's identity travels with it.

\section{Apparatus Accountability}
\label{se:apparatus-accountability}

When multiple records are combined, the apparatus must say which record is
responsible for which claim. Without that accountability, any reading can be
explained after the fact by any convenient story.

LIGO is the large-scale example. Two four-kilometer arms form an interferometer.
A gravitational wave is not recorded as a direct photograph of spacetime. It is
recorded as a differential change in arm length, extracted from an apparatus
that also carries seismic vibration, thermal noise, laser shot noise, mirror
motion, environmental monitors, timing systems, and calibration lines. The
apparatus grammar says which disturbances can count as signal and which must be
classified as noise or veto.

That grammar is not optional. If a truck drives past the facility, the ground
motion may enter sensors. If a mirror actuator glitches, the interferometer may
respond. If a laser component drifts, the record changes. The experiment must
carry enough auxiliary records to hold the main record accountable. Otherwise a
strain reading could be attributed to anything.

The same accountability exists in ordinary instruments. A speedometer depends
on tire radius, sensor state, clock state, and filtering logic. If the display
is wrong, one asks which part of the apparatus owns the failure. Is the tone
wheel damaged? Is the sensor missing pulses? Is the tire size wrong? Is the ECU
using an old calibration? Accountability is the ability to localize failure
without collapsing the whole experiment into mystery.

This section's lesson is that an experimental state has internal roles. Signal,
noise, calibration, veto, carrier, and display are not interchangeable. They are
parts of a grammar. The state is meaningful because the grammar says which
records may modify which claims.

Apparatus accountability also protects discovery. If every anomaly can be
dismissed as apparatus, nothing can be learned. If no anomaly can be assigned to
apparatus, every glitch becomes a new law. A good experiment lives between
these errors. It carries enough structure to say: this discrepancy belongs to
the sensor; this one belongs to the environment; this one survives the
instrument checks and therefore demands physical interpretation.

\section{Compatible Records}
\label{se:compatible-records}

Records become powerful when their constraints close. The speedometer,
clock, and odometer of \S\ref{se:experimental-state} are compatible when
they agree under the relation distance = speed times time. If the
speedometer reads 60 mph for one hour and the odometer advances only 40
miles, the state is incompatible. The contradiction does not reveal a new
kind of speed; it reveals a problem in the records or their conditions ---
miscalibration, failure, a wrong interval, wheel slip on a dynamometer,
readings from different trips. The grammar forces a diagnosis rather than
admitting a new physical fact.

Chemistry gives a sharper example because compatibility appears as ratio.
Hydrogen and oxygen combine as water in a fixed proportion. If an experiment
claims to have produced water while carrying incompatible mass ratios, the
records do not close. The apparatus may be leaking, the sample contaminated, or
the product misidentified. The chemical record has a structural constraint.

\begin{phenom}{Stoichiometry Effect}
\label{ph:stoichiometry-effect}

\PhStatement Chemical records compose only when their ratios are compatible
with the reaction structure. The product identity constrains the admissible
combination of reactant records.

\PhOrigin Stoichiometry emerged from the laws of definite and multiple
proportions: chemical compounds combine in fixed ratios by mass and count.

\PhObservation A laboratory record claiming water formation must reconcile the
hydrogen and oxygen inputs with the product. Excess reactant, missing mass, or
unexpected byproducts are not decorative details; they are compatibility checks.

\PhConstraint The experiment may not collapse all masses into one total and
call the result a chemical state. Source identities and reaction roles must be
preserved.

\PhConsequence Composition is structured. Multiple physical records become one
experimental state only when their ratios, roles, and sources fit the grammar
of the process being measured.

\end{phenom}

The Stoichiometry Effect shows why compatibility is not merely agreement in a
loose sense. It is agreement under a process grammar. Hydrogen has one role,
oxygen another, product another, residue another. A single total mass would
erase the structure the experiment needs.

The same is true of the car. Speed, distance, time, tire radius, and pulse count
must keep their roles. If all are collapsed into a single ``trip quality''
number, the state can no longer diagnose failure. The grammar has been lost.

\section{Unresolved States Held Open}
\label{se:unresolved-states-held-open}

Some experimental states are compatible precisely because they have not yet
been forced to choose.

The double-slit experiment is the canonical example. Before which-path
detection, the apparatus is arranged so that both paths contribute to the final
interference pattern. Treating the particle as if it had already chosen one
slit destroys the structure the experiment is designed to carry. The unresolved
state is not ignorance that should be patched over. It is a physical condition
of the experiment.

This is difficult because ordinary record-keeping wants closure. We want to
ask: which slit? But if the apparatus has not made that distinction admissible,
the question is not part of the record. A detector can be added. Then the
which-path record becomes admissible, and the interference pattern changes. The
new apparatus carries a different experimental state.

The speedometer has a quieter version. During a filtering window, the exact
rate may be unresolved. The instrument carries a pending state until enough
pulses arrive to update the display. Closing too early produces jitter or false
precision. Holding open too long produces lag. The instrument's design is a
policy for when an unresolved state becomes a reading.

An unresolved state is therefore not a failure of grammar. It is one of the
things grammar must support. The apparatus must be able to say: this value is
not yet closed; this path distinction is not admissible; this placeholder is
constrained but unresolved. Premature closure is as much an error as
incompatibility.

This prepares later quantum chapters, but it is already present in ordinary
measurement. A pending lab result, a provisional calibration, a filtered speed
estimate, and an undetected path all require the same discipline: carry the
unresolved state without pretending it has been resolved.

\section{Source / Compile / Execute}
\label{se:source-compile-execute}

An experiment also has stages. The protocol is not the instrument, and the
instrument is not the run.

The protocol is the source. It says how the experiment is to be performed:
which sample, which instrument, which calibration, which timing, which
exclusion rules, which analysis. The calibrated apparatus is the compiled form:
the protocol instantiated in hardware and procedure. The actual measurement run
is execution: the moment when the apparatus interacts with the world and
produces records.

Confusing these stages creates errors. A beautiful protocol does not guarantee
a clean run. A calibrated instrument does not guarantee that the sample was
prepared correctly. A successful run does not prove that the protocol was
general enough for future cases. Source, compile, and execute are related, but
they are distinct records.

The speedometer again helps. The source is the design and calibration rule:
count pulses, scale by circumference, divide by time. The compiled instrument
is the installed sensor, ECU program, tire calibration, and display. The
execution is a particular drive. If the tire is replaced, the compiled
instrument changes even if the source rule does not. If the car skids, the
execution may violate assumptions the compiled instrument normally relies on.

The grammar of experiments therefore preserves not only record types but stages
of realization. This is why lab notebooks separate protocol, apparatus,
calibration, raw data, and analysis. They are not redundant descriptions of one
thing. They are the source identities that let the experiment be checked.

Chapter 4 can now close. Multiple records combine by preserving roles, sources,
stages, and unresolved states. They do not become a heap. They become an
experimental grammar.

\begin{bridge}

The chapter's question was how multiple records combine without collapsing into
one undifferentiated record.

The answer is that records combine as an experimental state only through a
grammar that preserves their source identities, roles, calibrations, and stages.
A weather station combines readings by time and site. A speedometer, odometer,
and clock combine through distance-rate-time consistency. LIGO combines signal,
noise, veto, and calibration records through apparatus accountability. A
double-slit experiment holds unresolved alternatives open until the apparatus
admits a which-path record.

With experimental grammar in place, the next question can be asked. Records now
combine, but they still carry finite distinctions. Chapter 5 asks how a finite
count of distinctions becomes an admissible notion of distance.

\end{bridge}

\begin{coda}{A Hospital Operating Room Checklist}

Before the incision, the operating room pauses.

The patient is on the table. The anesthesiologist has induced sleep. The surgeon
is scrubbed. Instruments are arranged. Imaging is on the screen. The room is
ready to act, but action is not yet admissible.

The checklist begins.

Patient identity is confirmed. Procedure is named. Surgical site is marked.
Allergies are read. Antibiotics are checked. Imaging is confirmed. Instrument
counts are recorded. Anesthesia status is reported. Each item is a distinct
record. None is allowed to dissolve into the others. A correct name does not
substitute for a correct site mark. A complete instrument count does not
substitute for imaging. A prepared surgeon does not substitute for a verified
patient identity.

The checklist is not bureaucracy added to medicine from the outside. It is the
grammar that lets the procedure become one admissible state. The surgery can
begin only when the records are compatible.

If the site mark disagrees with the consent form, the disagreement is not a new
medical discovery. It is an incompatibility. If the instrument count is missing,
the state is unresolved. If the imaging belongs to another patient, the carrier
has failed. If the antibiotic time is outside the protocol window, the
calibration of the procedure is wrong. Each failure has a source.

The operating room therefore performs the same structure as the experimental
state. The protocol is the source. The prepared room is the compiled apparatus.
The incision is execution. The checklist preserves the distinction among these
stages so that action does not begin from a heap of good intentions.

Unresolved states are allowed, but they are held open. If the implant size is
not yet confirmed, the team does not pretend it is known. If a lab value is
pending, the procedure may wait or proceed under a named exception. The
placeholder is carried as a placeholder. It is not silently promoted into a
completed record.

The checklist also protects accountability. If a sponge is missing at closing,
the count record owns that failure. If the wrong image was displayed, the
imaging record owns it. If an allergy was missed, the pre-operative record owns
it. Because the records did not collapse into one another, the team can
diagnose the error.

The operating room is not a weather station, a speedometer, or a particle
experiment. But it shows why composition requires grammar. A safe procedure is
not many correct marks thrown together. It is a compatible state whose records
preserve their source identities until action is admissible.

Only then does the surgeon begin.

\end{coda}
